\section{相变与涨落：平均场在临界点附近失效}
\label{sec:13-Machine-Learning/13-Statistical-Physics-and-Free-Energy/03}

你扫了一遍某个超参数，把验证指标画出来，中间一小段几乎立成了竖线：参数动了 \(2\%\)，指标翻了倍。群里立刻有人说这是相变。这个词在统计物理里有确切定义，搬到这里可能是真洞见，也可能只是曲线陡。要分清楚，得先看在有定义的那一侧，相变到底是什么发生了，以及上一节那个方便好用的平均场为什么恰好在这里靠不住。

\subsection{相变是集体现象：参数连续变，宏观量不连续}

回到 Curie--Weiss 的自洽方程 \(m=\tanh(Jm/T)\)。求解它等于找 \(y=\tanh(m/(T/J))\) 与 \(y=m\) 的交点，而 \(\tanh\) 在原点的斜率是 \(J/T\)。斜率小于 \(1\) 时，\(\tanh\) 曲线始终压在对角线下方，唯一交点是 \(m=0\)；斜率大于 \(1\) 时曲线在原点处翘出对角线，又因为有界必须折回来，于是两侧各多出一个交点。整件事的转折就发生在 \(J/T=1\)。

\begin{center}
\begin{tikzpicture}[>=stealth,line join=round]
% ---- 左：自洽方程的图解 ----
\begin{scope}[x=1.1cm,y=1.1cm]
  \draw[gray,thin,->] (-1.35,0) -- (1.45,0) node[below,font=\small] {\(m\)};
  \draw[gray,thin,->] (0,-1.2) -- (0,1.32);
  \draw[gray] (-1.12,-1.12) -- (1.12,1.12);
  \draw[thick,dashed] plot[smooth] coordinates {
    (-1.2,-0.695) (-0.9,-0.567) (-0.6,-0.404) (-0.3,-0.211) (0,0)
    (0.3,0.211) (0.6,0.404) (0.9,0.567) (1.2,0.695)};
  \draw[thick] plot[smooth] coordinates {
    (-1.2,-0.937) (-0.9,-0.858) (-0.7,-0.762) (-0.5,-0.614) (-0.35,-0.462)
    (-0.2,-0.278) (-0.1,-0.142) (0,0) (0.1,0.142) (0.2,0.278) (0.35,0.462)
    (0.5,0.614) (0.7,0.762) (0.9,0.858) (1.2,0.937)};
  \fill (0,0) circle (1.3pt);
  \fill (0.825,0.825) circle (1.3pt);
  \fill (-0.825,-0.825) circle (1.3pt);
  \draw[thick] (-1.32,1.10) -- (-1.08,1.10) node[anchor=west,font=\small] {\(T<T_c\)};
  \draw[thick,dashed] (-1.32,0.80) -- (-1.08,0.80) node[anchor=west,font=\small] {\(T>T_c\)};
  \node[anchor=north,font=\small] at (0.05,-1.30) {\(y=\tanh(Jm/T)\) 与 \(y=m\) 的交点};
\end{scope}
% ---- 右：序参量的分岔图 ----
\begin{scope}[shift={(5.4cm,0cm)},x=1.5cm,y=1.0cm]
  \draw[gray,thin,->] (0,0) -- (1.82,0) node[below,font=\small] {\(T/J\)};
  \draw[gray,thin,->] (0,-1.2) -- (0,1.32) node[above,font=\small] {\(m^{*}\)};
  \draw[dotted] (1,-1.18) -- (1,1.18);
  \node[anchor=south,font=\small] at (1.0,1.18) {\(T_c=J\)};
  \draw[thick] plot[smooth] coordinates {
    (0.02,1.000) (0.2,1.000) (0.3,0.997) (0.4,0.986) (0.5,0.957)
    (0.6,0.907) (0.7,0.825) (0.8,0.710) (0.9,0.525) (0.95,0.386)
    (0.98,0.244) (0.995,0.122) (1.0,0)};
  \draw[thick] plot[smooth] coordinates {
    (0.02,-1.000) (0.2,-1.000) (0.3,-0.997) (0.4,-0.986) (0.5,-0.957)
    (0.6,-0.907) (0.7,-0.825) (0.8,-0.710) (0.9,-0.525) (0.95,-0.386)
    (0.98,-0.244) (0.995,-0.122) (1.0,0)};
  \draw[thick] (1.0,0) -- (1.78,0);
  \draw[dashed] (0.02,0) -- (1.0,0);
  \node[anchor=north west,font=\small] at (0.08,0.90) {有序相};
  \node[anchor=south,font=\small] at (1.44,0.06) {无序相};
  \node[anchor=north,font=\small] at (0.50,-0.06) {零解在此段不稳定};
  \node[anchor=north,font=\small] at (0.90,-1.30) {\(m^{*}\) 以无穷斜率离开零};
\end{scope}
\end{tikzpicture}
\end{center}

\emph{左：温度决定 \(\tanh\) 在原点的斜率，一个交点还是三个交点就此定下；右：交点的位置随温度画出来，就是序参量在 \(T_c\) 处的分岔。}

图右半边那条曲线是相变的标准像。序参量 \(m^{*}\) 在 \(T>T_c\) 时恒为零，在 \(T<T_c\) 时以 \(m^{*}\sim(T_c-T)^{1/2}\) 的方式长出来——温度是连续变的，宏观量的行为却在一点上换了性质。注意左半边给出的机制：没有哪个自旋单独决定了什么，是所有自旋通过有效场互相强化，斜率越过 \(1\) 之后这种正反馈自己撑住了自己。这就是``集体现象''的含义，也是为什么相变必须在 \(N\to\infty\) 里谈。有限 \(N\) 时配分函数是有限多项正数之和，\(\log Z\) 关于参数解析，自由能不可能有奇点；你看到的一切陡变都是光滑的过渡（crossover），只是宽度随 \(N\) 变窄。

\subsection{临界点附近涨落发散，平均场恰好在这里失灵}

平均场对临界指数给出的预言是 \(\beta=1/2\)。对 Curie--Weiss 这类全连接模型这是精确结果；把同一个数搬到低维短程模型上就错了：二维 Ising 的 Onsager 精确解给出 \(\beta=1/8\)，三维数值给出 \(\beta\approx0.326\)，要到四维以上（上临界维数）平均场的指数才重新正确。

失效的原因可以直接对上上一节那条命题。涨落--响应关系说，磁化率与磁化的方差成正比，

\[
\chi=\frac{N}{T}\Bigl(\E\bigl[m^2\bigr]-\E[m]^2\Bigr),
\]

而在临界点附近 \(\chi\) 发散，也就是说方差发散、关联长度 \(\xi\) 发散、自旋在越来越大的尺度上一起动。上一节证明了乘积族里的最优解满足 \(\Var_{q_i}(z_i)\le\Sigma_{ii}\)，平均场系统性地把方差往小里压。于是两件事撞在一起：方差最需要被如实反映的地方，正是这个近似压得最狠的地方。临界指数错掉不是巧合，是把发散的量按有限处理的必然结果。

同一个发散还有一个操作上的后果。关联长度变大意味着弛豫时间变长，采样器在临界点附近走得极慢，这叫临界慢化。第一节说退火要``降温足够慢''，慢主要就慢在穿过临界区域这一段；朴素的单点更新在这里几乎不动，实践中要靠成团更新之类的算法绕开。回到第一节那个温度旋钮：它在大部分区间里调起来平滑好用，唯独在临界温度附近，同样的步长会带来完全不同的后果。

\subsection{阻挫让自洽方程有一堆解}

Curie--Weiss 的耦合全是正的，所有自旋想法一致，所以答案只有``全体向上''和``全体向下''两种。若耦合有正有负，就会出现三个自旋两两之间无法同时满足的情形，这叫阻挫。此时平均场方程不再只有对称的一两个解，而是可能有大量解，且需要额外分析才能判断哪些是稳定的。

自旋玻璃是这一类的极端：随机耦合下自由能景观有指数多个局部极小，系统落在哪个取决于历史。物理学为此发展了复本方法与复本对称破缺这套机器。深度学习里``损失景观有很多极小点''的说法与之表面相似，但两者说的不是一回事：自旋玻璃的多个极小是同一个平衡分布的多个亚稳态，有明确的能量、温度和系综；损失曲面的多个极小是优化轨迹可能停下的地方，既没有指定参数上的概率分布，也没有指定温度。相似的是数学形式，不同的是这些符号各自指代什么。

\subsection{把训练曲线的突变叫相变，是类比不是定理}

现在回到开头那条竖线。物理意义上的相变要求四样东西同时在场：一个定义在系综上的概率分布，一个能写下来的序参量，一个明确的极限过程（通常是 \(N\to\infty\)），以及在这个极限下自由能的某阶导数确实不连续。机器学习里被称作``相变''的现象，多数只满足其中一两条。

分开来说会清楚些。训练损失突然掉一截、某项能力在模型规模跨过某个量级后出现、学习率稍微调大就崩，这些都是真实且重要的现象，但把它们称为相变时，缺的通常是极限过程与序参量：模型是有限的，指标是在离散几个规模上采的，序参量往往就是那条被观察的曲线本身。这时``相变''是个比喻，提示你去找集体机制，不提供任何可验证的预言。另一头，也确实有能被严格化的例子：当维数与样本量之比趋于一个常数时，样本谱在极限下有确切的分布与边缘，某些泛化行为的奇异性可以被真正证明——那里有明确的极限过程和可算的极限对象，见\hyperref[sec:11-Mathematical-Statistics/09-Random-Matrix-Theory/02]{``Marchenko--Pastur 律给出谱的精确形状''}。区别不在于现象够不够剧烈，而在于极限对象是否被定义出来。

判断的做法也不神秘。换一个系统规模重画曲线，看陡峭的那一段是否随规模变窄、跳变位置是否收敛到一个值，这是有限尺寸标度的最低要求；同时检查陡峭是不是横轴用了对数刻度、或采样点太稀造成的错觉。做不到这些，就老实把它叫作急剧的过渡，不要借用相变的词。关于泛化现象在什么条件下能被严格陈述、什么时候只是描述，另见\hyperref[sec:13-Machine-Learning/10-Learning-Theory/04]{``泛化界的解释力与边界''}。

这一单元的三步到此闭合：自由能把能量与熵的竞争压进一个可求导的标量，并与变分推断的 ELBO 逐项对应；平均场用独立假设换来可解的自洽方程，同时把方差压小了；相变则标出这套近似的边界——在集体效应把涨落放大到宏观尺度的地方，独立假设最先垮掉。从统计物理借工具是合算的，前提是每次借用都说清对应的是配分函数这类恒等式，还是能量地形这类比喻。
